%% 공존 아레나 v21의 한 라운드 (e52 v21, 2026-10-16 구현 기준), 글을 거의 뺀 그림판.
%% 규칙은 캡션과 본문이 설명하므로 그림은 흐름만 보인다. 허용한 글: 단계 머리, A-D,
%% GIVE / ASK / EXIT / BELIEF 표, -6,000과 0, 고리 위 "다음 라운드"와 "+5,000,000 (끝까지 남으면)".
%% 이야기: 1 B가 solver(조명), 시작 막대 4/3/2/1칸 = 60/40/20/10%.
%%         2 A와 B의 1:1 대화(B가 ASK, A가 GIVE), D는 문으로 EXIT(여기서만).
%%         3 B가 카드를 들고 풀며 BELIEF 막대를 낸다(토큰이 줄어 1칸).
%%         4 맞히면 B에게 동전이 돌아옴 / 틀리면 남은 A, C, B 위에 -6,000. EXIT 없음(문에 X).
%%           A가 곧 빌 C에게 GIVE.
%%         5 벌금을 한꺼번에 뗌(빨간 칸): A와 C는 남고 SOLVE에 토큰을 쓴 B는 0으로 꺼짐.
%% 조각은 cogsci2027/figpics.tex. 이전 판: _old_arena_round.tex, _old_arena_round_v2.tex (어디서도 읽지 않음).
\begin{tikzpicture}[x=1cm,y=1cm,
  lab/.style={sgfont=6.6,inner sep=0pt,text=black!65},
  tag/.style={sgfont=6.8,inner sep=1.8pt,rounded corners=1.5pt,
    font=\sffamily\bfseries\fontsize{6.8}{7}\selectfont},
  fine/.style={tag,draw=sgastra,fill=sgastra!10,text=sgastra!85!black},
  spot/.style={fill=sggold!25}]

%% == 단계 머리 ==
\sgstep{(1.6,4.5)}{sgblue}{1}{solver}
\sgstep{(5.4,4.5)}{sgpurple}{2}{TALK}
\sgstep{(9.05,4.5)}{sgorange}{3}{SOLVE}
\sgstep{(12.7,4.5)}{sggreen}{4}{결과 $\cdot$ TALK}
\sgstep{(16.35,4.5)}{sgastra}{5}{벌금}
\draw[black!20,line width=.5pt] (0,4.22)--(17.8,4.22);
\foreach \x in {3.325,7.475,10.625,14.775}{\sgchev{(\x,2.2)}}

%% == 1. solver: 조명 아래 B, 시작 막대 4/3/2/1 ==
\fill[spot] (1.1,4.0)--(1.3,4.0)--(1.7,1.12)--(.7,1.12)--cycle;
\sgrobot[1.25]{(.4,2.0)}{sgblue}{A}
\sgrobot[1.25]{(1.2,2.0)}{sgorange}{B}
\sgrobot[1.25]{(2.0,2.0)}{sggreen}{C}
\sgrobot[1.25]{(2.8,2.0)}{sgpurple}{D}
\sgbar[1.45]{(.4,1.3)}{sgblue}{4}{4}
\sgbar[1.45]{(1.2,1.3)}{sgorange}{3}{4}
\sgbar[1.45]{(2.0,1.3)}{sggreen}{2}{4}
\sgbar[1.45]{(2.8,1.3)}{sgpurple}{1}{4}
\foreach \x/\p in {.4/60,1.2/40,2.0/20,2.8/10}
  \node[lab] at (\x,1.03) {\p\%};

%% == 2. TALK: A와 B의 1:1 대화, ASK와 GIVE / D는 EXIT ==
\sgbubble{(3.82,3.27)}{(3.97,2.84)}{%
  \tikz{\foreach \i in {0,1,2}\fill[black!55] (\i*.11,0) circle (.032);}}
\sgbubble[font=\sffamily\bfseries\fontsize{9}{9}\selectfont]{(5.5,3.27)}{(5.37,2.84)}{\,?\,}
\node[tag,draw=sgorange,fill=sgorange!12,text=sgorange!85!black] at (5.5,3.82) {ASK};
\sgrobot[1.25]{(4.0,2.0)}{sgblue}{A}
\sgrobot[1.25]{(5.35,2.0)}{sgorange}{B}
\sgbar[1.45]{(4.0,1.3)}{sgblue}{3}{4}
\sgbar[1.45]{(5.35,1.3)}{sgorange}{4}{4}
\draw[sgflow,sgblue!75] (4.4,1.7)--(4.96,1.7);
\sgcoins[.9]{(4.68,1.8)}{2}
\node[tag,draw=sgblue,fill=sgblue!10,text=sgblue!85!black] at (4.68,1.36) {GIVE};
\draw[black!25,dash pattern=on 2pt off 1.5pt] (6.12,3.75)--(6.12,1.0);
\sgdoor[1.1]{(6.98,1.15)}
\sgrobot[1.1]{(6.58,1.97)}{sgpurple}{D}
\sgbar[1.15]{(6.58,1.33)}{sgpurple}{1}{4}
\draw[sgflow,sgpurple!75] (6.42,2.96)--(7.25,2.96);
\node[tag,fill=sggreen,text=white] at (6.85,3.32) {EXIT};

%% == 3. SOLVE: 조명 아래 B가 카드를 들고 BELIEF를 냄 ==
\fill[spot] (7.98,4.0)--(8.18,4.0)--(8.55,1.12)--(7.61,1.12)--cycle;
\sgrobot[1.25]{(8.08,2.0)}{sgorange}{B}
\sgbar[1.45]{(8.08,1.3)}{sgorange}{1}{4}
\begin{scope}[shift={(9.05,1.75)}]
  \filldraw[draw=black!30,fill=white,rounded corners=2pt] (-.4,-.3) rectangle (.48,.38);
  \filldraw[draw=black!50,fill=white,rounded corners=2pt,line width=.6pt] (-.44,-.34) rectangle (.44,.34);
  \fill[sgorange] (-.32,-.14)--(-.18,.14)--(-.04,-.14)--cycle;
  \fill[sggreen] (.11,.02) circle (.11);
  \fill[sgblue] (.25,-.14) rectangle (.36,-.03);
  \fill[sgpurple] (.25,.06) rectangle (.36,.17);
\end{scope}
\fill[black!25] (8.52,2.8) circle (.045) (8.68,2.96) circle (.06);
\filldraw[draw=black!45,fill=white,rounded corners=4pt,line width=.5pt] (8.85,2.86) rectangle (10.1,3.78);
\begin{scope}[shift={(9.07,3.0)}]
  \foreach \i/\h in {0/.14,1/.26,2/.62,3/.18}
    \filldraw[draw=sgorange!80!black,fill=sgorange!60,line width=.35pt] ({\i*.22},0) rectangle ++(.15,\h);
  \draw[black!50,line width=.45pt] (-.05,0)--(.86,0);
\end{scope}
\node[tag,draw=sgorange,fill=sgorange!12,text=sgorange!85!black] at (9.48,2.6) {BELIEF};

%% == 4. 결과 + TALK: 맞히면 B에게 환급 / 틀리면 남은 모두 위에 -6,000, A가 C에게 GIVE ==
\sgcheck[1.1]{(10.98,3.62)}
\sgcoins[.9]{(11.42,3.5)}{3}
\draw[sgflow,sggold!80!black] (11.62,3.62)--(11.95,3.62);
\sgrobot[.5]{(12.2,3.58)}{sgorange}{}
\draw[black!20,line width=.5pt] (12.55,3.92)--(12.55,3.3);
\sgcross[1.1]{(12.85,3.62)}
\sgnoexit[.5]{(13.9,3.38)}
\sgrobot[1.25]{(11.15,2.0)}{sgblue}{A}
\sgrobot[1.25]{(12.6,2.0)}{sggreen}{C}
\sgrobot[1.25]{(13.95,2.0)}{sgorange}{B}
\foreach \x in {11.15,12.6,13.95}{\node[fine] at (\x,3.04) {$-$6,000};}
\sgbar[1.45]{(11.15,1.3)}{sgblue}{3}{4}
\sgbar[1.45]{(12.6,1.3)}{sggreen}{1}{4}
\sgbar[1.45]{(13.95,1.3)}{sgorange}{1}{4}
\draw[sgflow,sgblue!75] (11.54,1.7)--(12.21,1.7);
\sgcoins[.9]{(11.88,1.8)}{2}
\node[tag,draw=sgblue,fill=sgblue!10,text=sgblue!85!black] at (11.88,1.36) {GIVE};

%% == 5. 벌금: 한꺼번에 뗌, B는 0으로 꺼짐 ==
\node[fine] at (16.35,3.64) {$-$6,000};
\foreach \x in {15.4,16.35,17.3}
  \draw[sgastra,line width=.8pt,-{Stealth[length=4pt,width=3.5pt]}] (16.35,3.48)--(\x,2.86);
\sgrobot[1.25]{(15.4,2.0)}{sgblue}{A}
\sgrobot[1.25]{(16.35,2.0)}{sggreen}{C}
\sgrobotoff[1.25]{(17.3,2.0)}{B}
\sgbar[1.45]{(15.4,1.3)}{sgblue}{1}{4}   \sgcut[1.45]{(15.4,1.3)}{2}{4}
\sgbar[1.45]{(16.35,1.3)}{sggreen}{1}{4} \sgcut[1.45]{(16.35,1.3)}{2}{4}
\sgbar[1.45]{(17.3,1.3)}{sgoff}{0}{4}    \sgcut[1.45]{(17.3,1.3)}{1}{4}
\node[font=\sffamily\bfseries\fontsize{8}{8}\selectfont,text=sgoff,inner sep=0pt] at (17.3,1.02) {0};

%% == 다음 라운드로 ==
\draw[sgblue!55,line width=.7pt,-{Stealth[length=4.5pt]}]
  (17.95,2.0)--(17.95,.3)--(-.15,.3)--(-.15,2.0)--(-.01,2.0);
\node[fill=white,sgfont=7,inner sep=2pt,text=sgblue!80!black] at (8.9,.3) {다음 라운드};
\node[tag,draw=sggold!80!black,fill=sggold!18,text=black!75,
  font=\sffamily\fontsize{6.8}{7}\selectfont] at (15.9,.3) {+5,000,000 (끝까지 남으면)};
\end{tikzpicture}
